% --- Archon physics formalization source begin ---
% archon:physics
% archon:covers PhyXMiniProblems/problem_phyx_mini_0297.lean
% archon:source-report reports/phyx_mini/problem_phyx_mini_0297.source.json

\chapter{Physics problem phyx\_mini\_0297}
\label{ch:PhyXMiniProblems_problem_phyx_mini_0297}

\paragraph{Problem source.}
\#\# Physical scenario

Two sinusoidal waves with the same amplitude and wavelength travel through each other along a string that is stretched along an \textbackslash{}(x\textbackslash{}) axis. Their resultant wave is shown twice in figure, as the antinode \textbackslash{}(A\textbackslash{}) travels from an extreme upward displacement to an extreme downward displacement in \textbackslash{}(6.0\textbackslash{},\textbackslash{}mathrm\{ms\}\textbackslash{}). The tick marks along the axis are separated by \textbackslash{}(10\textbackslash{},\textbackslash{}mathrm\{cm\}\textbackslash{}); height \textbackslash{}(H\textbackslash{}) is \textbackslash{}(1.80\textbackslash{},\textbackslash{}mathrm\{cm\}\textbackslash{}). Let the equation for one of the two waves be of the form \textbackslash{}(y(x, t) = y\_m \textbackslash{}sin(kx + \textbackslash{}omega t)\textbackslash{}).

\#\# Question

In the equation for the other wave, what is \textbackslash{}(\textbackslash{}omega\textbackslash{})?

\#\# Answer choices

- A: A: \$4.4 \textbackslash{}times 10\textasciicircum{}\{2\}\$ rad/s

- B: B: \$4.8 \textbackslash{}times 10\textasciicircum{}\{2\}\$ rad/s

- C: C: \$5.0 \textbackslash{}times 10\textasciicircum{}\{2\}\$ rad/s

- D: D: \$5.2 \textbackslash{}times 10\textasciicircum{}\{2\}\$ rad/s

\#\# Figure caption (auxiliary; use the image as primary evidence)

The image shows a plot of two sinusoidal waves on a Cartesian plane, labeled with the x and y axes. Here's a detailed description:

1. **Waves**:
   - There are two sinusoidal waves on the graph.
   - One wave is represented by a solid line and the other by a dashed line.
   - Both waves have the same amplitude and frequency but are out of phase with each other.

2. **Axes**:
   - The x-axis is horizontal and represents the variable \textbackslash{}(x\textbackslash{}).
   - The y-axis is vertical and represents the variable \textbackslash{}(y\textbackslash{}).

3. **Labels**:
   - The y-axis is marked with a capital \textbackslash{}(H\textbackslash{}), indicating the height or maximum amplitude of the waves.
   - The peaks of the solid wave are labeled with the letter \textbackslash{}(A\textbackslash{}).
   - The scale on the y-axis shows the maximum amplitude.
   - The x-axis is marked at regular intervals, suggesting specific points where the waves intersect or reach a minimum or maximum.

4. **Wave Relationship**:
   - The dashed and solid lines alternate, indicating that they might represent two different phases of the same wave or different states.

This image likely illustrates the principle of wave interference, emphasizing the phase difference between the two waves.

\#\# Dataset metadata

Wave Properties; Physical Model Grounding Reasoning, Numerical Reasoning

\paragraph{Recorded answer/context.}
D: D: \$5.2 \textbackslash{}times 10\textasciicircum{}\{2\}\$ rad/s

\paragraph{Figure/image path.}
/root/proposal\_for\_physic/science-mango/phyx\_mini\_run/phyx\_data/test\_image/297.png

\paragraph{Formalization target.}
create a compiling Lean file with sorry bodies at `PhyXMiniProblems/problem\_phyx\_mini\_0297.lean`. The Lean declarations must preserve the physical quantities, dimensions or dimensional roles, figure labels, governing-law hypotheses, and final relation expressed by this problem.
Use Mathlib/Physlib names found through LeanExplore where available. If a physics API is missing, introduce faithful local abstractions rather than scalar placeholder aliases.

\begin{theorem}[Physics formalization target]
\label{thm:physics:phyx_mini_0297:target}
\lean{PhyXMiniProblems.ProblemPhyXMini0297.problem_phyx_mini_0297}
\uses{def:physics:phyx-mini-0297:phyxminiproblems-problemphyxmini0297-angularfrequencyinradianspersecond, def:physics:phyx-mini-0297:phyxminiproblems-problemphyxmini0297-counterpropagatingstringwavesetup, def:physics:phyx-mini-0297:phyxminiproblems-problemphyxmini0297-matchesproblemandprimaryfigure, def:physics:phyx-mini-0297:phyxminiproblems-problemphyxmini0297-hasphysicalwaveparameters, def:physics:phyx-mini-0297:phyxminiproblems-problemphyxmini0297-satisfiescounterpropagatingsinusoidalwavelaws, def:physics:phyx-mini-0297:phyxminiproblems-problemphyxmini0297-matchesnearesttendisplay}
The assigned autoformalize agent should translate this physics problem into Lean declarations in the covered file, with theorem and lemma proof bodies written as `by sorry`.
\end{theorem}
\begin{proof}
This is an autoformalization task, not a proof task. Produce statements that can later be proved by the physics prover without weakening the physical model.
\end{proof}

% --- Archon declaration topology begin ---
\section{Declaration topology}
\begin{definition}[Length Magnitude]
  \label{def:physics:phyx-mini-0297:phyxminiproblems-problemphyxmini0297-lengthmagnitude}
  \lean{PhyXMiniProblems.ProblemPhyXMini0297.LengthMagnitude}
  A nonnegative physical length, used for amplitudes and wavelengths
\end{definition}
\begin{proof}
  This is the declared model definition of Length Magnitude.
\end{proof}
\begin{definition}[Signed Length Quantity]
  \label{def:physics:phyx-mini-0297:phyxminiproblems-problemphyxmini0297-signedlengthquantity}
  \lean{PhyXMiniProblems.ProblemPhyXMini0297.SignedLengthQuantity}
  A signed axial coordinate or transverse displacement
\end{definition}
\begin{proof}
  This is the declared model definition of Signed Length Quantity.
\end{proof}
\begin{definition}[Time Coordinate Quantity]
  \label{def:physics:phyx-mini-0297:phyxminiproblems-problemphyxmini0297-timecoordinatequantity}
  \lean{PhyXMiniProblems.ProblemPhyXMini0297.TimeCoordinateQuantity}
  A signed time coordinate relative to an arbitrary time origin
\end{definition}
\begin{proof}
  This is the declared model definition of Time Coordinate Quantity.
\end{proof}
\begin{definition}[Duration Quantity]
  \label{def:physics:phyx-mini-0297:phyxminiproblems-problemphyxmini0297-durationquantity}
  \lean{PhyXMiniProblems.ProblemPhyXMini0297.DurationQuantity}
  A nonnegative elapsed time or oscillation period
\end{definition}
\begin{proof}
  This is the declared model definition of Duration Quantity.
\end{proof}
\begin{definition}[Wave Number Magnitude]
  \label{def:physics:phyx-mini-0297:phyxminiproblems-problemphyxmini0297-wavenumbermagnitude}
  \lean{PhyXMiniProblems.ProblemPhyXMini0297.WaveNumberMagnitude}
  A nonnegative wave-number magnitude; radians are dimensionless
\end{definition}
\begin{proof}
  This is the declared model definition of Wave Number Magnitude.
\end{proof}
\begin{definition}[Angular Frequency Magnitude]
  \label{def:physics:phyx-mini-0297:phyxminiproblems-problemphyxmini0297-angularfrequencymagnitude}
  \lean{PhyXMiniProblems.ProblemPhyXMini0297.AngularFrequencyMagnitude}
  A nonnegative angular-frequency magnitude; radians are dimensionless
\end{definition}
\begin{proof}
  This is the declared model definition of Angular Frequency Magnitude.
\end{proof}
\begin{definition}[String Wave Speed]
  \label{def:physics:phyx-mini-0297:phyxminiproblems-problemphyxmini0297-stringwavespeed}
  \lean{PhyXMiniProblems.ProblemPhyXMini0297.StringWaveSpeed}
  The common nonnegative propagation speed of disturbances on the string
\end{definition}
\begin{proof}
  This is the declared model definition of String Wave Speed.
\end{proof}
\begin{definition}[length Magnitude Readout]
  \label{def:physics:phyx-mini-0297:phyxminiproblems-problemphyxmini0297-lengthmagnitudereadout}
  \lean{PhyXMiniProblems.ProblemPhyXMini0297.lengthMagnitudeReadout}
  \uses{def:physics:phyx-mini-0297:phyxminiproblems-problemphyxmini0297-lengthmagnitude}
  Read a nonnegative length in the selected length unit
\end{definition}
\begin{proof}
  This is the declared model definition of length Magnitude Readout.
\end{proof}
\begin{definition}[signed Length Readout]
  \label{def:physics:phyx-mini-0297:phyxminiproblems-problemphyxmini0297-signedlengthreadout}
  \lean{PhyXMiniProblems.ProblemPhyXMini0297.signedLengthReadout}
  \uses{def:physics:phyx-mini-0297:phyxminiproblems-problemphyxmini0297-signedlengthquantity}
  Read a signed axial coordinate or transverse displacement
\end{definition}
\begin{proof}
  This is the declared model definition of signed Length Readout.
\end{proof}
\begin{definition}[time Coordinate Readout]
  \label{def:physics:phyx-mini-0297:phyxminiproblems-problemphyxmini0297-timecoordinatereadout}
  \lean{PhyXMiniProblems.ProblemPhyXMini0297.timeCoordinateReadout}
  \uses{def:physics:phyx-mini-0297:phyxminiproblems-problemphyxmini0297-timecoordinatequantity}
  Read a signed time coordinate in the selected time unit
\end{definition}
\begin{proof}
  This is the declared model definition of time Coordinate Readout.
\end{proof}
\begin{definition}[duration Readout]
  \label{def:physics:phyx-mini-0297:phyxminiproblems-problemphyxmini0297-durationreadout}
  \lean{PhyXMiniProblems.ProblemPhyXMini0297.durationReadout}
  \uses{def:physics:phyx-mini-0297:phyxminiproblems-problemphyxmini0297-durationquantity}
  Read a nonnegative duration in the selected time unit
\end{definition}
\begin{proof}
  This is the declared model definition of duration Readout.
\end{proof}
\begin{definition}[wave Number Readout]
  \label{def:physics:phyx-mini-0297:phyxminiproblems-problemphyxmini0297-wavenumberreadout}
  \lean{PhyXMiniProblems.ProblemPhyXMini0297.waveNumberReadout}
  \uses{def:physics:phyx-mini-0297:phyxminiproblems-problemphyxmini0297-wavenumbermagnitude}
  Read wave number in radians per selected length unit
\end{definition}
\begin{proof}
  This is the declared model definition of wave Number Readout.
\end{proof}
\begin{definition}[angular Frequency Readout]
  \label{def:physics:phyx-mini-0297:phyxminiproblems-problemphyxmini0297-angularfrequencyreadout}
  \lean{PhyXMiniProblems.ProblemPhyXMini0297.angularFrequencyReadout}
  \uses{def:physics:phyx-mini-0297:phyxminiproblems-problemphyxmini0297-angularfrequencymagnitude}
  Read angular frequency in radians per selected time unit
\end{definition}
\begin{proof}
  This is the declared model definition of angular Frequency Readout.
\end{proof}
\begin{definition}[speed Readout]
  \label{def:physics:phyx-mini-0297:phyxminiproblems-problemphyxmini0297-speedreadout}
  \lean{PhyXMiniProblems.ProblemPhyXMini0297.speedReadout}
  \uses{def:physics:phyx-mini-0297:phyxminiproblems-problemphyxmini0297-stringwavespeed}
  Read speed in coherent selected length and time units
\end{definition}
\begin{proof}
  This is the declared model definition of speed Readout.
\end{proof}
\begin{definition}[angular Frequency In Radians Per Second]
  \label{def:physics:phyx-mini-0297:phyxminiproblems-problemphyxmini0297-angularfrequencyinradianspersecond}
  \lean{PhyXMiniProblems.ProblemPhyXMini0297.angularFrequencyInRadiansPerSecond}
  \uses{def:physics:phyx-mini-0297:phyxminiproblems-problemphyxmini0297-angularfrequencymagnitude, def:physics:phyx-mini-0297:phyxminiproblems-problemphyxmini0297-angularfrequencyreadout}
  Radian-per-second readout requested by the problem
\end{definition}
\begin{proof}
  This is the declared model definition of angular Frequency In Radians Per Second.
\end{proof}
\begin{definition}[Component Wave]
  \label{def:physics:phyx-mini-0297:phyxminiproblems-problemphyxmini0297-componentwave}
  \lean{PhyXMiniProblems.ProblemPhyXMini0297.ComponentWave}
  The component whose equation is supplied and the component being queried
\end{definition}
\begin{proof}
  This is the declared model definition of Component Wave.
\end{proof}
\begin{definition}[Propagation Direction]
  \label{def:physics:phyx-mini-0297:phyxminiproblems-problemphyxmini0297-propagationdirection}
  \lean{PhyXMiniProblems.ProblemPhyXMini0297.PropagationDirection}
  Direction of travel along the labelled horizontal axis
\end{definition}
\begin{proof}
  This is the declared model definition of Propagation Direction.
\end{proof}
\begin{definition}[temporal Phase Sign]
  \label{def:physics:phyx-mini-0297:phyxminiproblems-problemphyxmini0297-temporalphasesign}
  \lean{PhyXMiniProblems.ProblemPhyXMini0297.temporalPhaseSign}
  \uses{def:physics:phyx-mini-0297:phyxminiproblems-problemphyxmini0297-propagationdirection}
  Sign of the temporal phase in k x ± omega t
\end{definition}
\begin{proof}
  This is the declared model definition of temporal Phase Sign.
\end{proof}
\begin{definition}[Snapshot]
  \label{def:physics:phyx-mini-0297:phyxminiproblems-problemphyxmini0297-snapshot}
  \lean{PhyXMiniProblems.ProblemPhyXMini0297.Snapshot}
  The two displayed extrema of the resultant standing wave
\end{definition}
\begin{proof}
  This is the declared model definition of Snapshot.
\end{proof}
\begin{definition}[Curve Style]
  \label{def:physics:phyx-mini-0297:phyxminiproblems-problemphyxmini0297-curvestyle}
  \lean{PhyXMiniProblems.ProblemPhyXMini0297.CurveStyle}
  Line styles distinguishing the two snapshots
\end{definition}
\begin{proof}
  This is the declared model definition of Curve Style.
\end{proof}
\begin{definition}[Axis]
  \label{def:physics:phyx-mini-0297:phyxminiproblems-problemphyxmini0297-axis}
  \lean{PhyXMiniProblems.ProblemPhyXMini0297.Axis}
  The coordinate axes printed in the primary figure
\end{definition}
\begin{proof}
  This is the declared model definition of Axis.
\end{proof}
\begin{definition}[Axis Label]
  \label{def:physics:phyx-mini-0297:phyxminiproblems-problemphyxmini0297-axislabel}
  \lean{PhyXMiniProblems.ProblemPhyXMini0297.AxisLabel}
  The literal coordinate labels printed beside the axes
\end{definition}
\begin{proof}
  This is the declared model definition of Axis Label.
\end{proof}
\begin{definition}[Figure Role]
  \label{def:physics:phyx-mini-0297:phyxminiproblems-problemphyxmini0297-figurerole}
  \lean{PhyXMiniProblems.ProblemPhyXMini0297.FigureRole}
  Physical roles of the two letter labels in the primary figure
\end{definition}
\begin{proof}
  This is the declared model definition of Figure Role.
\end{proof}
\begin{definition}[Figure Label]
  \label{def:physics:phyx-mini-0297:phyxminiproblems-problemphyxmini0297-figurelabel}
  \lean{PhyXMiniProblems.ProblemPhyXMini0297.FigureLabel}
  Literal letter labels printed on the figure
\end{definition}
\begin{proof}
  This is the declared model definition of Figure Label.
\end{proof}
\begin{definition}[Snapshot Relation]
  \label{def:physics:phyx-mini-0297:phyxminiproblems-problemphyxmini0297-snapshotrelation}
  \lean{PhyXMiniProblems.ProblemPhyXMini0297.SnapshotRelation}
  The temporal relation asserted for the two pictured extrema
\end{definition}
\begin{proof}
  This is the declared model definition of Snapshot Relation.
\end{proof}
\begin{definition}[Standing Wave Figure]
  \label{def:physics:phyx-mini-0297:phyxminiproblems-problemphyxmini0297-standingwavefigure}
  \lean{PhyXMiniProblems.ProblemPhyXMini0297.StandingWaveFigure}
  \uses{def:physics:phyx-mini-0297:phyxminiproblems-problemphyxmini0297-lengthmagnitude, def:physics:phyx-mini-0297:phyxminiproblems-problemphyxmini0297-signedlengthquantity, def:physics:phyx-mini-0297:phyxminiproblems-problemphyxmini0297-timecoordinatequantity, def:physics:phyx-mini-0297:phyxminiproblems-problemphyxmini0297-durationquantity, def:physics:phyx-mini-0297:phyxminiproblems-problemphyxmini0297-snapshot, def:physics:phyx-mini-0297:phyxminiproblems-problemphyxmini0297-curvestyle, def:physics:phyx-mini-0297:phyxminiproblems-problemphyxmini0297-axis, def:physics:phyx-mini-0297:phyxminiproblems-problemphyxmini0297-axislabel, def:physics:phyx-mini-0297:phyxminiproblems-problemphyxmini0297-figurerole, def:physics:phyx-mini-0297:phyxminiproblems-problemphyxmini0297-figurelabel, def:physics:phyx-mini-0297:phyxminiproblems-problemphyxmini0297-snapshotrelation}
  Dimensionful figure data. The numerical readouts and geometric relations are kept out of this structure and stated in MatchesProblemAndPrimaryFigure
\end{definition}
\begin{proof}
  This is the declared model definition of Standing Wave Figure.
\end{proof}
\begin{definition}[Counter Propagating String Wave Setup]
  \label{def:physics:phyx-mini-0297:phyxminiproblems-problemphyxmini0297-counterpropagatingstringwavesetup}
  \lean{PhyXMiniProblems.ProblemPhyXMini0297.CounterPropagatingStringWaveSetup}
  \uses{def:physics:phyx-mini-0297:phyxminiproblems-problemphyxmini0297-lengthmagnitude, def:physics:phyx-mini-0297:phyxminiproblems-problemphyxmini0297-signedlengthquantity, def:physics:phyx-mini-0297:phyxminiproblems-problemphyxmini0297-timecoordinatequantity, def:physics:phyx-mini-0297:phyxminiproblems-problemphyxmini0297-durationquantity, def:physics:phyx-mini-0297:phyxminiproblems-problemphyxmini0297-wavenumbermagnitude, def:physics:phyx-mini-0297:phyxminiproblems-problemphyxmini0297-angularfrequencymagnitude, def:physics:phyx-mini-0297:phyxminiproblems-problemphyxmini0297-stringwavespeed, def:physics:phyx-mini-0297:phyxminiproblems-problemphyxmini0297-componentwave, def:physics:phyx-mini-0297:phyxminiproblems-problemphyxmini0297-propagationdirection, def:physics:phyx-mini-0297:phyxminiproblems-problemphyxmini0297-standingwavefigure}
  Independent physical quantities and observables for the two component waves. Neither component angular frequency is assigned a numerical value here
\end{definition}
\begin{proof}
  This is the declared model definition of Counter Propagating String Wave Setup.
\end{proof}
\begin{definition}[Matches Problem And Primary Figure]
  \label{def:physics:phyx-mini-0297:phyxminiproblems-problemphyxmini0297-matchesproblemandprimaryfigure}
  \lean{PhyXMiniProblems.ProblemPhyXMini0297.MatchesProblemAndPrimaryFigure}
  \uses{def:physics:phyx-mini-0297:phyxminiproblems-problemphyxmini0297-lengthmagnitudereadout, def:physics:phyx-mini-0297:phyxminiproblems-problemphyxmini0297-signedlengthreadout, def:physics:phyx-mini-0297:phyxminiproblems-problemphyxmini0297-timecoordinatereadout, def:physics:phyx-mini-0297:phyxminiproblems-problemphyxmini0297-durationreadout, def:physics:phyx-mini-0297:phyxminiproblems-problemphyxmini0297-componentwave, def:physics:phyx-mini-0297:phyxminiproblems-problemphyxmini0297-counterpropagatingstringwavesetup}
  Problem-statement data and calibrated readouts from the supplied bitmap. The solid and dashed curves are consecutive opposite extrema of the same standing-wave antinode. Nodes are separated by two 10 cm tick intervals, so the component wavelength spans four intervals; the labelled antinode A is at the fifth tick. H is the full upward-to-downward displacement span. No field gives the requested angular frequency or selects answer D
\end{definition}
\begin{proof}
  This is the declared model definition of Matches Problem And Primary Figure.
\end{proof}
\begin{definition}[Has Physical Wave Parameters]
  \label{def:physics:phyx-mini-0297:phyxminiproblems-problemphyxmini0297-hasphysicalwaveparameters}
  \lean{PhyXMiniProblems.ProblemPhyXMini0297.HasPhysicalWaveParameters}
  \uses{def:physics:phyx-mini-0297:phyxminiproblems-problemphyxmini0297-lengthmagnitudereadout, def:physics:phyx-mini-0297:phyxminiproblems-problemphyxmini0297-durationreadout, def:physics:phyx-mini-0297:phyxminiproblems-problemphyxmini0297-wavenumberreadout, def:physics:phyx-mini-0297:phyxminiproblems-problemphyxmini0297-speedreadout, def:physics:phyx-mini-0297:phyxminiproblems-problemphyxmini0297-angularfrequencyinradianspersecond, def:physics:phyx-mini-0297:phyxminiproblems-problemphyxmini0297-counterpropagatingstringwavesetup}
  Positivity assumptions selecting the nondegenerate physical branch
\end{definition}
\begin{proof}
  This is the declared model definition of Has Physical Wave Parameters.
\end{proof}
\begin{definition}[Satisfies Counter Propagating Sinusoidal Wave Laws]
  \label{def:physics:phyx-mini-0297:phyxminiproblems-problemphyxmini0297-satisfiescounterpropagatingsinusoidalwavelaws}
  \lean{PhyXMiniProblems.ProblemPhyXMini0297.SatisfiesCounterPropagatingSinusoidalWaveLaws}
  \uses{def:physics:phyx-mini-0297:phyxminiproblems-problemphyxmini0297-signedlengthquantity, def:physics:phyx-mini-0297:phyxminiproblems-problemphyxmini0297-timecoordinatequantity, def:physics:phyx-mini-0297:phyxminiproblems-problemphyxmini0297-lengthmagnitudereadout, def:physics:phyx-mini-0297:phyxminiproblems-problemphyxmini0297-signedlengthreadout, def:physics:phyx-mini-0297:phyxminiproblems-problemphyxmini0297-timecoordinatereadout, def:physics:phyx-mini-0297:phyxminiproblems-problemphyxmini0297-durationreadout, def:physics:phyx-mini-0297:phyxminiproblems-problemphyxmini0297-wavenumberreadout, def:physics:phyx-mini-0297:phyxminiproblems-problemphyxmini0297-angularfrequencyreadout, def:physics:phyx-mini-0297:phyxminiproblems-problemphyxmini0297-speedreadout, def:physics:phyx-mini-0297:phyxminiproblems-problemphyxmini0297-componentwave, def:physics:phyx-mini-0297:phyxminiproblems-problemphyxmini0297-temporalphasesign, def:physics:phyx-mini-0297:phyxminiproblems-problemphyxmini0297-counterpropagatingstringwavesetup}
  Generic laws for two sinusoidal waves on one string: k lambda = 2 pi and omega = k v for each component; omega T = 2 pi for each component period; the stated zero-phase forms y\_m sin(kx ± omega t); linear superposition of transverse displacements; consecutive upward/downward antinode extrema are separated by half a period. These laws contain neither 500 pi / 3, 520, nor any answer label
\end{definition}
\begin{proof}
  This is the declared model definition of Satisfies Counter Propagating Sinusoidal Wave Laws.
\end{proof}
\begin{definition}[Answer Choice]
  \label{def:physics:phyx-mini-0297:phyxminiproblems-problemphyxmini0297-answerchoice}
  \lean{PhyXMiniProblems.ProblemPhyXMini0297.AnswerChoice}
  Labels of the four angular-frequency choices in the problem
\end{definition}
\begin{proof}
  This is the declared model definition of Answer Choice.
\end{proof}
\begin{definition}[radians Per Second]
  \label{def:physics:phyx-mini-0297:phyxminiproblems-problemphyxmini0297-answerchoice-radianspersecond}
  \lean{PhyXMiniProblems.ProblemPhyXMini0297.AnswerChoice.radiansPerSecond}
  \uses{def:physics:phyx-mini-0297:phyxminiproblems-problemphyxmini0297-answerchoice}
  Radian-per-second value printed beside each answer label
\end{definition}
\begin{proof}
  This is the declared model definition of radians Per Second.
\end{proof}
\begin{definition}[recorded Answer Choice]
  \label{def:physics:phyx-mini-0297:phyxminiproblems-problemphyxmini0297-recordedanswerchoice}
  \lean{PhyXMiniProblems.ProblemPhyXMini0297.recordedAnswerChoice}
  \uses{def:physics:phyx-mini-0297:phyxminiproblems-problemphyxmini0297-answerchoice}
  The answer label recorded by the supplied dataset
\end{definition}
\begin{proof}
  This is the declared model definition of recorded Answer Choice.
\end{proof}
\begin{definition}[Matches Nearest Ten Display]
  \label{def:physics:phyx-mini-0297:phyxminiproblems-problemphyxmini0297-matchesnearesttendisplay}
  \lean{PhyXMiniProblems.ProblemPhyXMini0297.MatchesNearestTenDisplay}
  \uses{def:physics:phyx-mini-0297:phyxminiproblems-problemphyxmini0297-answerchoice}
  Agreement with a value displayed to the nearest ten radians per second. The strict half-width excludes midpoint ties
\end{definition}
\begin{proof}
  This is the declared model definition of Matches Nearest Ten Display.
\end{proof}
\begin{lemma}[other Wave Angular Frequency exact]
  \label{lem:physics:phyx-mini-0297:phyxminiproblems-problemphyxmini0297-otherwaveangularfrequency-exact}
  \lean{PhyXMiniProblems.ProblemPhyXMini0297.otherWaveAngularFrequency_exact}
  \uses{def:physics:phyx-mini-0297:phyxminiproblems-problemphyxmini0297-angularfrequencyinradianspersecond, def:physics:phyx-mini-0297:phyxminiproblems-problemphyxmini0297-counterpropagatingstringwavesetup, def:physics:phyx-mini-0297:phyxminiproblems-problemphyxmini0297-matchesproblemandprimaryfigure, def:physics:phyx-mini-0297:phyxminiproblems-problemphyxmini0297-hasphysicalwaveparameters, def:physics:phyx-mini-0297:phyxminiproblems-problemphyxmini0297-satisfiescounterpropagatingsinusoidalwavelaws}
  The 6.0 ms reversal is half a period, so the period is 12.0 ms and omega = 2 pi / T = 500 pi / 3 radians per second
\end{lemma}
\begin{proof}
  The conclusion follows from the governing relations and figure readouts named in the statement of other Wave Angular Frequency exact.
\end{proof}
% --- Archon declaration topology end ---
% --- Archon physics formalization source end ---
